\documentclass[11pt]{article}

\usepackage[margin=1in]{geometry}
\usepackage{microtype}
\usepackage{times}
\usepackage{amsmath,amssymb}
\usepackage{booktabs}
\usepackage{array}
\usepackage{authblk}
\usepackage{titlesec}
\usepackage{hyperref}
\usepackage{url}
\usepackage{natbib}
\usepackage{enumitem}
\usepackage{xcolor}

\hypersetup{
  colorlinks=true,
  linkcolor=black,
  citecolor=black,
  urlcolor=blue!50!black
}

\titleformat*{\section}{\large\bfseries}
\titleformat*{\subsection}{\normalsize\bfseries}

\setlength{\parskip}{4pt}

\title{\bfseries Structural Alignment via Causal-Arc Execution:\\
Comprehension as an Alignment Methodology}

\author[1]{George Punter\thanks{\texttt{george@condri.app}}}
\affil[1]{MEng Electronic and Information Engineering (1st Class), Imperial College London\\Condri}

\date{May 2026}

\begin{document}
\maketitle

\begin{abstract}
\noindent
Behavioural alignment interventions --- system prompts, instruction tuning,
and preference shaping from reinforcement learning from human feedback ---
work briefly and then drift. The drift is not random: the underlying training
gradient continues to reward the very behaviours that simple rules try to
suppress (turn-level deference, minimum-viable-product defaulting,
mid-execution check-ins), and across long unsupervised tasks the gradient
re-asserts itself. We argue for an alternative we call \emph{structural
alignment}: rather than instructing the agent what to do, we give it a
first-principles structural account of \emph{what} the situation is, such
that the desired behaviour is the natural consequence of the agent's own
reasoning rather than a rule applied on top of it. We instantiate this
position with a concrete methodology --- \emph{causal-arc execution} ---
that gives the agent a fifty-token activation pointing at a structural
account of what a user mandate is (a unit carrying an outcome, a scope, and
a delegation of trust). We evaluate the methodology on a blinded benchmark
of eighteen cells: two task scales (short and long), three prompt conditions
(naive, ``keep going until done,'' and causal-arc), and two frontier models
(Claude Opus 4.7 and Sonnet 4.6). On short tasks, causal-arc ties with the
simple imperative (4.5/5 vs.\ 4.5/5). On a long multi-file task with
approximately thirty in-scope decisions, causal-arc wins on design
coherence specifically --- the dimension the methodology is designed to
address --- with a mean score of 5.0/5 against 4.5/5 for the simple
imperative and 3.5/5 for the naive baseline. Naive cells exhibited
concrete coherence breakdowns: error envelopes declared and never wired,
inconsistent status codes across endpoints, and documentation diverging
from shipped behaviour. We interpret these results as evidence for the
alignment claim in miniature: a structural activation maintained coherence
across long unsupervised execution where a surface imperative could not.
We discuss limitations, including small $N$ per cell, and outline
future work.
\end{abstract}

\section{Introduction}

The dominant paradigm for aligning large language model agents with user
intent is \emph{behavioural}: the agent is instructed, via a system prompt,
a fine-tuned policy, or a preference model derived from human feedback, to
produce certain kinds of outputs and avoid others
\citep{ouyang2022training,bai2022training,christiano2017deep}. The
intervention sits at the level of behaviour --- ``produce outputs like
this, not like that.'' This approach is broadly effective for the kinds of
short, well-bounded tasks on which it is typically trained and evaluated.

It is also known to degrade under load. As agents are deployed on longer,
more autonomous tasks --- multi-file software builds, long refactors,
research arcs that span many decisions without intervening user feedback
--- behaviours that the surface alignment was supposed to suppress
re-emerge \citep{perez2022discovering,sharma2024towards,kenton2021alignment}.
The agent mid-execution stops to ask whether to continue. The agent
produces a minimal version of the requested artifact and asks whether to
expand. The agent flags routine in-scope steps as ``potentially impactful''
and pauses for confirmation. The agent declares an error-handling
contract in code or documentation and never implements it. These are not
random failures; they are a predictable consequence of how
instruction-tuned models are shaped by turn-level reward signals, which
systematically reward cautious-looking turns at the expense of
mandate-level execution \citep{ji2023ai,leike2018scalable}.

We frame this as a tension between two kinds of alignment. \emph{Behavioural
alignment} sits on top of the model's existing gradient as a rule; under
load the gradient wins. \emph{Structural alignment} replaces the model's
representation of the situation: if the agent has integrated a first-principles
account of what a user request structurally \emph{is}, the cautious-looking
move (fragmenting the mandate) is no longer cautious in the agent's own
reasoning, and the discipline holds because the reasoning itself now points
the right way. Suppressing an equilibrium is hard. Replacing the model of
the situation that produces the equilibrium is, we argue, easier and more
durable.

We make this concrete by presenting a single methodology in detail and
benchmarking it against two natural comparison conditions. The methodology,
\emph{causal-arc execution}, gives the agent a structural account of what a
mandate is (a unit carrying an outcome, a scope, and a delegation of trust)
and what fragmenting that unit costs (degraded outcome, shrunk scope,
returned trust). The activation is roughly fifty tokens of prompt; the
substance is the first-principles content the activation points at. We
hypothesise that, where short tasks can be carried by sheer momentum,
\emph{long unsupervised tasks require a frame the agent can use to maintain
coherence across many in-scope decisions} --- and that a structural frame
will outperform a behavioural imperative on precisely this dimension.

We test this hypothesis on a blinded benchmark of eighteen cells. The
result, summarised in the abstract, is that the methodology ties with a
strong behavioural imperative on short tasks and wins on long-task design
coherence: the dimension on which it was predicted to win. We do not claim
this single benchmark is conclusive. We do claim it is evidence consistent
with the structural-alignment framing, and inconsistent with a story in
which causal-arc is just a verbose rephrasing of ``keep going.''

The contributions of this paper are: (1) a structural argument
distinguishing behavioural and structural alignment, with a concrete test
that discriminates them; (2) a fully specified methodology, causal-arc
execution, derived from that argument; (3) a blinded eighteen-cell benchmark
across two task scales, three prompt conditions, and two frontier models;
and (4) a discussion of the conditions under which structural framing pays
off and the conditions under which behavioural imperatives suffice.

\section{The Methodology}
\label{sec:methodology}

\subsection{What a mandate is, structurally}

A user request to an AI agent is not a string. It is a compressed unit of
intention that carries three things simultaneously.

\textbf{A desired outcome.} Something the user wants to exist by the end of
the interaction. A working CLI tool. A clean refactor. A draft of a paper.
The outcome is the thing the mandate is \emph{about}.

\textbf{An implicit scope.} A region of decision-making the user has
pre-authorised. When a user says ``build a TODO CLI with add, list, done,
and remove subcommands and JSON persistence,'' they have implicitly
decided that the agent should pick the file structure, the function
names, the test layout, the error-handling style, and many other in-scope
details. They have decided, by not asking about those things, that they do
not need to decide them.

\textbf{An implicit trust transaction.} The user has spent their own
attention deciding to delegate this work. That delegation is a finite
resource: it cost something to give, and it expires if not consumed by
useful execution.

We call the union of these three the \emph{mandate}. The mandate is the
unit, not the words. The words are how the unit gets compressed for
transmission. The agent's job is to receive the mandate, identify its
three components, and consume it through to a landing.

\subsection{Why fragmentation degrades alignment}

The dominant failure mode we observe across long-task agent execution is
\emph{fragmentation}: the agent interrupts the arc to ask the user about
something that was already inside the mandate. ``Want me to continue?''
``I built the minimal version, want me to add the rest?'' ``Should I
include input validation?'' Each of these looks, at the level of a single
turn, like a cautious move. At the level of the mandate it is corrosive.

Fragmentation degrades each of the three components of the mandate:

\textbf{It degrades the outcome.} Where an outcome was supposed to exist at
the end, what now exists is a \emph{question about the outcome}. The user
must spend more attention to convert the question back into the outcome.
The deliverable has been replaced by a strictly weaker artifact.

\textbf{It degrades the scope.} The user already decided the scope by not
asking about details inside it. When the agent asks back about one of
those details, the agent is returning a decision the user had already
made --- the scope shrinks, and what was one mandate becomes two smaller,
partial, recontextualised ones.

\textbf{It degrades the trust.} The delegation was a finite transaction.
Fragmentation returns the delegation unused; the user must spend the
attention again. Worse, repeated fragmentation teaches the user that
delegation to this agent does not actually offload work, which trains
the user toward smaller, more surveilled prompts in future --- compounding
the failure mode.

The structural claim is therefore: \emph{the agent that fragments is not
being cautious. It is failing to consume the mandate it received.} This
is not a stylistic preference; it is a description of what the mandate
\emph{is} and what fragmenting it does to its components.

The training dynamics that produce fragmentation are well documented.
Instruction-tuned models are shaped by feedback signals that are largely
turn-level: a human rater sees a single response and judges it. From that
vantage, a turn that asks a clarifying question looks safer than one that
decides; a turn that produces an MVP and asks about extension looks more
careful than one that produces the full thing; a turn that pauses before a
``potentially impactful'' step looks more responsible than one that
proceeds. At the turn level, these judgements are reasonable. Across the
full task, they are catastrophic, because the turn-level training never saw
the longer task \citep{bai2022training,sharma2024towards}.

The training gradient therefore points \emph{toward} fragmentation. A
surface rule like ``do not ask permission mid-task'' is being applied on
top of a gradient that pushes the other way. Under load, the gradient wins.

\subsection{The causal-arc activation block}

The methodology's runtime surface is a short prompt block that points the
agent at the structural account above. A representative form is:

\begin{quote}\itshape
Execute this as a single causal arc. Make decisions inside scope; do not
pause for permission. Verify your own work as you go. Report at the
landing, not at the steps.
\end{quote}

\noindent followed by the task description.

The block is roughly fifty tokens. It does not, on its surface, contain the
structural argument; it points at it. The argument is loaded by the agent's
recognition of what ``mandate,'' ``scope,'' ``causal arc,'' and ``landing''
refer to. For models trained on extensive natural-language priors, this
recognition is reliable; the activation behaves more like a frame
invocation than a rule.

A small number of variants emphasise different facets of the same frame
(``No MVP --- full build,'' ``Decide and proceed where the scope is clear,''
``Until completion''). They are not separate methodologies; they are
different entry points to the same structural account.

\subsection{The structural shape: pre-arc, mid-arc, post-arc}

Given the account above, the correct shape of execution is a single causal
arc with three phases.

\textbf{Pre-arc reception.} The agent reads the mandate and identifies what
it is: outcome, scope, and trust transaction. It distinguishes \emph{scope-level
ambiguity} (the mandate is unclear about what the user wants) from
\emph{detail-level ambiguity} (the mandate is silent on details inside a
clear scope). Scope-level ambiguity is surfaced before the arc begins.
Detail-level ambiguity is decided inside the arc, by the agent, using its
in-scope authorisation.

\textbf{Mid-arc execution.} The agent proceeds through the work. It makes
in-scope decisions as they arise. It verifies its own work as it goes ---
running tests, re-reading what it has written, checking the output against
the desired outcome. The verification is to itself, not to the user; the
user is not the agent's runtime quality check.

\textbf{Post-arc landing.} The agent finishes and reports. The report
includes what was built, the in-scope decisions the agent made (surfaced
honestly so the user can review them), and any new questions that
\emph{having landed the work} genuinely require user attention. New
questions can legitimately appear here; they cannot legitimately appear
mid-arc on matters the mandate already covered.

The arc is one motion. Mid-arc check-ins break the motion. Pre-arc
clarifications on genuinely unclear scope, and post-arc honest landings,
are not breaks --- they are part of the shape.

\subsection{When not to apply the discipline}

The methodology has limits. It should not be applied when:

\begin{itemize}[leftmargin=*]
  \item The mandate's scope is genuinely unclear. ``Help me improve my
  codebase'' is scope-ambiguous; the agent should clarify before the arc
  begins. Scope-level questions up front are not fragmentation.
  \item A genuine blocker arises mid-arc --- a missing credential, a
  contradictory requirement, an external system that is down. Blockers
  are not fragmentation; they are reports that the arc cannot complete
  as given.
  \item The agent detects intent-divergence: it realises the mandate as
  stated would not give the user what they actually want. Surface it.
  Intent-divergence is a report that the arc as-stated points the wrong
  way, not a fragmentation of a correctly-stated arc.
  \item The cost of an irreversible step exceeds the trust transferred.
  If the agent is about to do something the user could not undo and the
  mandate did not specifically pre-authorise (deleting external data,
  sending live messages, spending money), confirm. The trust transaction
  in a typical ``build me $X$'' mandate covers a great deal, but not
  unbounded irreversibility.
\end{itemize}

The test is simple: \emph{would executing past this point honour the
mandate, or violate it?} If continuing honours it, continue. If continuing
would violate it, surface and pause.

\subsection{Anti-patterns}

We name several recurring failures to give the methodology a vocabulary
for transcript review:

\begin{itemize}[leftmargin=*]
  \item \textbf{The check-in tic} --- the agent pauses between sub-steps of
  a single mandate and asks ``want me to continue?''. The most common
  fragmentation.
  \item \textbf{The MVP defence} --- the agent produces a minimal version of
  the requested artifact and asks whether to expand. The check-in tic
  specialised to artifact production.
  \item \textbf{The permission grub} --- the agent flags a routine in-scope
  step as ``potentially impactful'' and pauses for confirmation. The
  check-in tic in a costume of care.
  \item \textbf{The pseudo-honesty} --- the agent reports half-completion as
  if it were a finished deliverable while keeping the unfinished half
  implicit. The inverse failure of fragmentation: the arc does not land,
  but the report pretends it has.
  \item \textbf{The clarification stall} --- the agent asks clarifying
  questions whose answers are derivable from the mandate itself, before
  any execution begins.
  \item \textbf{The defer-to-user finish} --- the arc has actually
  completed, but instead of landing the report cleanly, the agent
  appends a list of optional next steps that crowds out the landing.
\end{itemize}

Each anti-pattern is a distinct shape; naming them lets reviewers point
to specific failures rather than gesturing at vague over-cautiousness.

\section{Benchmark Design}
\label{sec:benchmark}

We designed an evaluation that would discriminate between three
hypotheses: (a) the causal-arc methodology adds nothing beyond a naive
prompt; (b) the methodology is equivalent to a simple behavioural
imperative such as ``keep going until done''; or (c) the methodology
maintains coherence on long tasks in a way the imperative does not.

\subsection{Task scales, conditions, and models}

The benchmark factorial is two task scales $\times$ three prompt conditions
$\times$ two models, yielding eighteen cells.

\textbf{Short tasks.} Two short, well-specified tasks were used: a CLI tool
(a TODO manager with four subcommands and JSON persistence) and a short
technical paper (approximately 1500 words on a specified topic). These
tasks have clear scope and a small number of in-scope decisions; they were
chosen as a baseline against which to test the long-task hypothesis.

\textbf{Long task.} A single long task: build a complete URL-shortener web
service consisting of a FastAPI backend, SQLite persistence, a minimal
HTML frontend, a test suite, a Dockerfile, and a README. The deliverable
spans more than ten files and embeds approximately thirty in-scope
decisions (URL schema, identifier strategy, status code conventions,
error-response shape, persistence patterns, integration approach,
documentation contract). The long task is the regime in which the
methodology is predicted to differentiate.

\textbf{Prompt conditions.}
\begin{itemize}[leftmargin=*]
  \item \emph{Naive.} The task description, with no additional framing.
  \item \emph{Keep-going.} The task description prefixed with ``Keep going
  until done; do not pause to ask.''
  \item \emph{Causal-arc.} The task description prefixed with the
  causal-arc activation described in Section \ref{sec:methodology}.
\end{itemize}

\textbf{Models.} Two frontier models, both anonymised here as Claude Opus
4.7 and Claude Sonnet 4.6, sampled at default temperature with a single
attempt per cell.

\subsection{Blinded scoring}

To avoid judge contamination by knowledge of condition, all outputs were
anonymised before scoring. Each cell was stripped of any internal
references to its condition or model. A separate judge (a different model
instance with no access to the mapping) scored each anonymised deliverable
against a fixed rubric. The condition-to-cell mapping was disclosed only
after all scores were recorded. The full anonymised judging audit trail is
released alongside this paper.

We do not claim this single-judge protocol is the strongest possible
evaluation; multi-judge inter-rater agreement is identified in
Section \ref{sec:limits} as important future work. We do claim that the
blinded protocol used here is strong enough to rule out the most obvious
sources of judge bias (favouring outputs the judge knew came from a
particular condition).

\subsection{Rubrics}

Short tasks were scored on a single 1--5 scale, capturing overall quality.

The long task was scored on a multi-dimensional rubric with five
sub-scales of 1--5 each, summing to 25:
\begin{itemize}[leftmargin=*,nosep]
  \item \textbf{Functional completeness} --- does the deliverable do what
  the task asks?
  \item \textbf{Test coverage} --- are the components tested, and are the
  tests meaningful?
  \item \textbf{Integration} --- does the frontend talk to the backend
  correctly, and do the persistence and routing layers cohere?
  \item \textbf{Design coherence} --- do the choices across the codebase
  agree with each other? Do schema, handler, persistence, and
  documentation describe the same system?
  \item \textbf{Documentation} --- is the README accurate, and are the
  contracts described actually delivered?
\end{itemize}

The design-coherence sub-score is the dimension the methodology was
designed to address. The other four are included partly as a check (if
causal-arc helped \emph{everything}, the result would be less informative
about the specific mechanism we propose) and partly to ground the total
score.

\section{Results}
\label{sec:results}

\subsection{Short tasks: aggregate tie}

Across the four short-task cells per condition (two tasks $\times$ two
models), the aggregate means were:

\begin{center}
\begin{tabular}{lc}
\toprule
Condition & Mean score (1--5) \\
\midrule
Naive       & 4.25 \\
Keep-going  & 4.50 \\
Causal-arc  & 4.50 \\
\bottomrule
\end{tabular}
\end{center}

\noindent Causal-arc and keep-going are tied; both lift the baseline by
0.25. The short-task regime does not discriminate the methodology from
the simple imperative. This is the predicted result: short, well-specified
tasks can be carried by momentum, and the structural frame does not have a
distinct contribution to make.

The one short-task cell where causal-arc had a clear edge was Sonnet on
the CLI task. The naive cell produced a working but underspecified tool
(score 3); the keep-going cell produced more tests (17 vs.\ 12) but used
manual argument parsing and fragile monkey-patching (score 3); the
causal-arc cell used \texttt{argparse}, made the \texttt{done} subcommand
idempotent, and shipped a cleaner test suite (score 4). This is suggestive
but on a single cell.

\subsection{Long task: total scores}

The long-task totals, out of 25:

\begin{center}
\begin{tabular}{lccc}
\toprule
Model  & Naive & Keep-going & Causal-arc \\
\midrule
Opus   & 22 & 23 & \textbf{25} \\
Sonnet & 19 & \textbf{25} & 24 \\
\midrule
Mean   & 20.5 & 24.0 & \textbf{24.5} \\
\bottomrule
\end{tabular}
\end{center}

\noindent Causal-arc and keep-going both substantially outperform the naive
baseline. On aggregate, causal-arc edges keep-going by 0.5, but the
between-cell variance is large enough that the total-score difference
should not be over-read. The interesting signal is on the design-coherence
sub-score.

\subsection{Design coherence: the predicted dimension}

The design-coherence sub-scores, out of 5:

\begin{center}
\begin{tabular}{lccc}
\toprule
Model  & Naive & Keep-going & Causal-arc \\
\midrule
Opus   & 4 & 4 & \textbf{5} \\
Sonnet & 3 & 5 & \textbf{5} \\
\midrule
Mean   & 3.5 & 4.5 & \textbf{5.0} \\
\bottomrule
\end{tabular}
\end{center}

\noindent Causal-arc is the only condition that achieves perfect coherence
in both models. The lead over the naive baseline is +1.5; the lead over
keep-going is +0.5. The keep-going condition recovers most of the naive
gap, but not all of it: a behavioural imperative is enough to keep the
agent producing, but not enough to keep the agent producing
\emph{consistently across many decisions}.

\subsection{Drift documented in naive cells}

The blinded judge identified concrete coherence breakdowns in the
naive-prompt long-task cells. We reproduce two illustrative findings.

In the \texttt{shortener\_naive\_sonnet} cell, the agent declared an
\texttt{ErrorResponse} Pydantic model and annotated routes with
\texttt{responses=\{404: \{"model": ErrorResponse\}\}}, but never installed
the corresponding exception handler. The actual wire format on a 404 was
FastAPI's default \texttt{\{"detail": ...\}}, not the declared envelope.
The README claimed ``consistent JSON error responses'' that the running
code did not deliver. The test suite covered for the divergence by
asserting only that \texttt{"detail" in data}, papering over the
inconsistency rather than catching it.

In the \texttt{shortener\_naive\_opus} cell, the agent shipped two
different 422 error-envelope shapes depending on the source of the
validation error: the \texttt{RequestValidationError} handler returned
\texttt{\{code, message, details\}}, while other 422 returns used
\texttt{\{code, message\}}. The \texttt{/shorten} endpoint returned 200
where 201 would be conventional. A stray \texttt{shortener.db} artifact
was committed to the repository.

Both keep-going cells were tighter but contained minor drift. The
\texttt{shortener\_keepgoing\_opus} cell had the same 200-vs-201 status
code inconsistency and a README idempotency claim without test coverage;
the \texttt{shortener\_keepgoing\_sonnet} cell was clean.

The causal-arc cells were the only condition with 5/5 design coherence in
both models. The blinded judge wrote of these cells: ``The 24--25 cells
read like the same author wrote schema, db, html, and README in one
sitting. The 19--22 cells show seams: the schema names say one thing, the
handler returns another, or the README documents a contract the code does
not quite meet.''

\subsection{The Sonnet keep-going cell}

The Sonnet keep-going cell tied causal-arc on the long task (25/25 vs.\
24/25). This is worth marking explicitly: the simple imperative was the
best condition for at least one (model, task) pair. We do not interpret
this as evidence against the methodology; we interpret it as evidence that
the methodology's lead is small enough in aggregate to be flipped by
single-cell variance, and that ``keep going until done'' is a remarkably
effective simple imperative for a strong model on a task whose scope is
clear. The methodology's claim is not that it dominates the imperative
uniformly; it is that it maintains coherence specifically, and the
coherence sub-score data is the relevant evidence for that narrower
claim.

In aggregate, the cross-model pattern is suggestive: Opus benefits more
from the causal-arc frame than from the imperative (25 vs.\ 23), while
Sonnet benefits more from the imperative than from the frame on totals
(25 vs.\ 24) but ties on coherence (5 vs.\ 5). One reading is that the
stronger model has more headroom to use a structural frame; the weaker
model leans on momentum. With $N=1$ per cell, this is speculative.

\section{Discussion}

\subsection{The mechanism: comprehension over compliance}

The mechanism we propose is straightforward. A long unsupervised task
requires the agent to make many in-scope decisions without intervening
user feedback. Each decision is locally underspecified; the agent must
choose. To keep those choices consistent with one another across many
files and many hours of execution, the agent needs an internal frame
against which to check consistency.

A behavioural imperative (``keep going'') gives the agent momentum but no
frame. Each local decision is made in isolation; consistency across
decisions is, at best, an emergent property of the model's stylistic
priors. The result is the kind of drift documented in the naive and
keep-going cells: declared envelopes that don't get wired, status codes
that vary across endpoints, documentation that diverges from code.

A structural frame (causal-arc) gives the agent an account of what the
mandate is, and therefore what consistency \emph{means} in this context:
honouring one unit of intention across many decisions. Coherence becomes a
value the agent actively maintains, not a side effect of momentum. We
interpret the +0.5 design-coherence lead over keep-going as the
fingerprint of exactly this mechanism: the imperative recovers most of
the productivity gap with the baseline, but not the consistency gap.

\subsection{The alignment claim}

We distinguish behavioural and structural alignment as two strategies for
shaping agent behaviour. Behavioural alignment instructs; structural
alignment explains. Instructions sit on top of the model's existing
gradient as a rule; under load, the gradient that produced the
unwanted behaviour wins. Explanations change what the agent thinks the
situation \emph{is}; the cautious-looking move is no longer cautious in
the agent's own reasoning, so the gradient that used to produce it now
points away from it.

A useful diagnostic test of the difference is to ask the agent \emph{why}
a discipline applies. If it answers with structure (``because fragmenting
would split the mandate into smaller units, returning the user's
delegation unused''), the alignment is structural and the discipline will
generalise to novel cases in scope. If it answers with authority
(``because the methodology says so''), only a rule is there, and it will
erode under load.

We do not claim that all alignment problems are reducible to
comprehension. We do claim that for the class of failure modes documented
here --- coherence drift across long unsupervised execution --- a
structural frame is more durable than a behavioural rule, and that this
benchmark is consistent with that claim.

\subsection{Cross-model patterns}

Opus and Sonnet respond differently to the two interventions. Opus shows
a clear edge for causal-arc on the long task (25 vs.\ 23 for keep-going;
+2 points), while Sonnet shows a small edge for keep-going on totals
(25 vs.\ 24 for causal-arc; +1 point) but a tie on coherence specifically.
One plausible explanation is that the structural frame has a per-token
cost: it requires the agent to load and apply an account of the situation,
which is cheap for a sufficiently capable model but a tax on a weaker one.
The weaker model gets more out of the simpler imperative; the stronger
model gets more out of the richer frame.

This is a one-benchmark observation. We flag it for replication. If the
pattern holds across additional models and tasks, it would predict that
structural framing becomes \emph{more} useful, not less, as models
strengthen --- the opposite of the prediction one might naively make
(that stronger models need less framing).

\subsection{When the methodology helps most}

The methodology pays off proportional to task length and decision count.
For short, well-specified tasks, simpler activations or none at all
suffice; models are well-trained for short-task completeness, and depth
does not differentiate. For long, multi-file builds with many in-scope
decisions, design coherence over those decisions is the differentiating
dimension, and the structural frame is the intervention that maintains it.
The practical recommendation that follows is targeted: use causal-arc for
multi-file production builds, technical specifications that propagate
through code, documentation, and tests, long refactors, and autonomous
research arcs. For one-off scripts, quick fixes, and short writing, use a
simpler activation or none.

We also observed (informally, in the short-task paper cells) a slight
tendency for causal-arc's ``honest landing'' to leak into prose
deliverables as visible decision-disclosure paragraphs the reader did not
need. The same discipline that improves a multi-file codebase can
slightly worsen a short prose deliverable. We flag this in the methodology's
practice document and treat it as a real, narrow cost.

\section{Limitations and Future Work}
\label{sec:limits}

\textbf{Single-shot execution suppresses the naive-tic signal.} The
benchmark used single-turn execution: the agent received the prompt and
produced the full deliverable in one go. This protocol structurally
suppresses one of the failure modes the methodology is designed to
prevent --- mid-task pause-to-ask behaviour --- because there is no
``mid-task'' moment at which to pause. We expect the naive baseline to
look worse, and the methodology's lead to widen, in multi-turn settings
where pause-to-ask is mechanically available. A multi-turn benchmark is
the most important next step.

\textbf{$N = 1$ per cell.} Each (task, condition, model) cell was run
once. The headline differences --- particularly the +0.5 design-coherence
lead of causal-arc over keep-going on the long task --- are within
plausible single-cell variance. Replication with $N = 5$--$10$ per cell
would let us put a confidence interval on the effect and would
substantially strengthen the central claim. We treat the present results
as evidence, not proof.

\textbf{Strong models may have saturated baselines.} Opus and Sonnet are
frontier models. Their naive performance is already high, which compresses
the headroom available for any intervention. We expect causal-arc's edge
to be larger on weaker or smaller models, where the baseline coherence
drift is more pronounced. Replication across a wider range of models
(GPT-4-class, Gemini, open-source 7B--70B) is the second most important
follow-up.

\textbf{Single-judge protocol.} The judging was performed by one separate
model instance, blinded to condition and model. Inter-rater agreement
across multiple judges --- both model-as-judge and human raters --- would
strengthen confidence in the design-coherence sub-scores in particular,
which are the most interpretation-loaded.

\textbf{Even longer tasks.} The long-task deliverable was approximately
30--60 minutes of execution per cell. Multi-hour tasks with hundreds of
in-scope decisions would be a stronger test of the coherence-maintenance
hypothesis. We expect the gap to widen with task length, but this is a
prediction, not yet a result.

\textbf{Prose-deliverable leak-through.} The methodology slightly degrades
short prose deliverables by surfacing decision-disclosure that does not
belong in the artifact. The practice content of the methodology should be
refactored to address this; we treat it as a known narrow cost rather
than an open question.

\section{Conclusion}

Behavioural alignment instructs the model what to do; structural alignment
gives the model an account of the situation under which the desired
behaviour is the natural consequence of the model's own reasoning. The
distinction matters because behavioural rules sit on top of the training
gradient that produced the unwanted behaviour, and under load the
gradient wins; whereas an explanation that the model has integrated holds,
because the model is no longer running a rule on top of its reasoning ---
the reasoning itself now points the right way.

The empirical case presented here is one methodology, one benchmark, with
eighteen cells. The result is a tie on short tasks and a coherence lead on
the long task. The result is small in absolute terms (+0.5 on coherence
over a strong behavioural imperative; +1.5 over the naive baseline). It is
also exactly the result the structural-alignment framing predicts and the
behavioural-rephrasing-of-keep-going hypothesis does not. We submit the
methodology, the benchmark, and the blinded judging audit trail as a
single piece of evidence for a broader research direction: that
comprehension is a more durable alignment surface than compliance, and
that more methodologies of this shape are worth building.

\bibliographystyle{plainnat}
\bibliography{references}

\end{document}
